\section{Discussion}
\label{sec:discussion}

\subsection{Why the wind changes from the station to the road}
The results separate three causes of the difference between synoptic and roadside wind. The measurement height and the local roughness explain a large part in a forested lowland such as western Sweden, where synoptic stations stand in open terrain and roads run through forest; the logarithmic wind profile and the blending-height concept \citep{Wieringa1986} capture this with fixed physical constants and are insensitive to their exact values. Second, each reference station has a persistent exposure of its own, of a factor up to two in either direction (\noSfMin--\noSfMax{} in Norway), which must be removed before stations can be combined; a single factor per station suffices. Third, and most important at the roadside, the wind is reduced by forest at two distinct scales: the trees next to the road, within a few tree heights, and the forest fraction of the surrounding landscape within about \SI{1}{km}. The first scale corresponds to the sheltered zone behind forest edges and windbreaks \citep{HeislerDeWalle1988,DupontBrunet2008}; the second to the roughness of the landscape that controls the near-surface wind through the internal boundary layer. The absence of any effect at intermediate radii of \SIrange{300}{500}{m}, found in both countries, supports this two-scale view.

In mountainous Norway the landscape forest fraction dominates, the physical roughness conversion adds little, and explicit terrain indices did not help. A plausible reason is that the forest fraction already acts as a proxy for terrain: in central Norway forest grows in sheltered valleys, while exposed plateaus and coasts are open. The near-tree coefficient changes sign between the countries for the same reason; it should be read together with the forest fraction, not as an independent physical effect.

\subsection{Combining stations}
The single largest improvement comes from combining all reference stations within about \SI{100}{km} rather than taking the nearest one, after the site factors have made them comparable. A single station, even nearby, carries its own exposure and its own local turbulence; averaging many corrected stations recovers the regional wind, which the roadside exposure model then localises. This agrees with the common practice of using the regional flow from an analysis or model rather than a single observation, and explains why the MET Nordic analysis, which combines observations with a numerical model, is a strong benchmark at vehicle height.

\subsection{Transferability}
The model form---the same predictors and the same physical steps---works in both landscapes, and its coefficients are stable in time within each country. The coefficients, however, are partly regional: transferred to the other country, the model keeps a clear advantage over the nearest station in Sweden but only a small one in Norway. For a new region with roadside stations the coefficients should therefore be refitted, which requires only open data; for a region without roadside stations, the coefficients of the most similar landscape should be used and the uncertainty reported. The feature distributions of the training stations define the domain of applicability: mountain passes, bridges and urban street canyons are underrepresented in both data sets.

\subsection{Use for vehicles}
Three conclusions follow for vehicle applications. First, the exposure features decorrelate along the road over about \SI{250}{m} in both countries, so segments for vehicle-based analyses should not be much longer than \SIrange{200}{300}{m} if local exposure is to be resolved; longer segments average over different environments. Second, the station-based estimate is clearly better than the nearest station but, at vehicle height and on the two roads tested, not better than the MET Nordic analysis, mainly because of large direction errors in a mountain valley. The two sources are complementary: the station-based model contributes local exposure, the analysis the regional flow and its direction. Third, the error of the static model has a regional component that changes with the weather (an out-of-time bias of 10--17\,\%). Both points suggest an online fusion, in which a Kalman filter combines the analysis, the corrected stations and roadside observations, and the exposure model of this paper maps the result to the road segment.

For the air drag of a tractor--semitrailer, the corrected wind halves the median drag error at the Swedish stations and reduces it by about a quarter in Norway. These numbers depend on the drag model; Model~1 of \citet{Askerdal2024crosswind} is their recommendation when no CFD or wind-tunnel data exist and is likely to overstate the crosswind sensitivity of a specific vehicle, so the relative improvement is more reliable than the absolute errors.

\subsection{Limitations}
The Swedish record covers only \seHours{} hours in November, because the Trafikverket API keeps seven days; the Swedish results confirm the Norwegian ones but do not cover seasons. Statens vegvesen sensor heights are taken as the nominal \SI{10}{m} reported in Frost; deviations would be absorbed by the intercept but could affect the roughness conversion. Roadside stations are not vehicle-height truth, and the vehicle validation covers only three days on two roads. The drag errors are model-based, not measured.
